\paragraph*{04.05.07.04 Finanzmanagement- und Berichtssystem für das Projekt entwickeln, einrichten und
aufrechterhalten}
\kcilabel{para:04.05.07.04_Finanzmanagement}{04.05.07.04}
\label{para:04.05.07.04_Finanzmanagement}

% Task
\par Die \gls{pl} und die \gls{tpl} waren für die finanzielle Steuerung des Projekts verantwortlich, während ich als \gls{pm} die operative Umsetzung sowie den stabilen Betrieb eines effektiven Finanzmanagement- und Berichtssystems sicherstellte. Meine Aufgabe bestand darin, ein System zu entwickeln und einzurichten, das die kontinuierliche Erfassung, Überwachung und Berichterstattung der Projektkosten ermöglicht und damit eine transparente, belastbare Grundlage für Steuerung und Entscheidungsfindung schafft.

% Action
\par Durch die Analyse der bestehenden Finanzprozesse und -strukturen im \gls{s4} identifizierte ich die Anforderungen an das Finanzmanagement- und Berichtssystem. In enger Zusammenarbeit mit der \gls{pl} und der \gls{tpl} übernahm ich das im \gls{inex} etablierte Kostenreporting auf \gls{epic}-Ebene als methodische Grundlage und passte es auf den Projektkontext an. Das Reporting umfasste einen wellenspezifischen Soll-Ist-Vergleich sowie einen Forecast im Burn-Down-Format. Hierfür strukturierte ich die Entwicklungsaufwände auf Basis der Anforderungen, nutzte die bestehende Anforderungszuordnung als Kostenzuordnung und verdichtete die Ergebnisse zu einer belastbaren Entscheidungsgrundlage. Zur Aufrechterhaltung des Systems definierte ich einen regelmäßigen Berichtszyklus, klare Verantwortlichkeiten für Datenlieferung und Qualitätssicherung sowie einen verbindlichen Übergang der Ergebnisse in die bestehenden Entscheidungs- und Berichtsstrukturen der \gls{pl}.

% Result
\par In den folgenden Monaten etablierte ich das Finanzmanagement- und Berichtssystem erfolgreich, indem ich die relevanten Datenquellen für die Kostenerfassung identifizierte, die notwendigen Schnittstellen zu den bestehenden Systemen einrichtete und die Berichtsprozesse verbindlich definierte. Das System ermöglichte eine kontinuierliche Überwachung der Projektkosten, einen transparenten Soll-Ist-Vergleich und eine belastbare Prognose der Kostenentwicklung je Welle. Durch den stabilen Regelbetrieb konnten Abweichungen früher erkannt, adressiert und in den zuständigen Gremien zeitnah entschieden werden. Dies erhöhte die Transparenz über die finanziellen Aspekte des Projekts und verbesserte die Qualität der finanziellen Steuerung.

% Reflect
\par Die Implementierung des Finanzmanagement- und Berichtssystems hat gezeigt, dass enge Zusammenarbeit zwischen den Projektbeteiligten und eine klare Definition der Anforderungen entscheidend für den Erfolg sind. Da dies nicht immer gegeben war, konnte ich die Einführung nicht so schnell vorantreiben, wie es für die erste Umsetzungswelle ideal gewesen wäre. Für künftige Vorhaben würde ich daher bereits in der Planungsphase verbindliche Anforderungen an Reporting, Datenqualität und Entscheidungswege festlegen, die notwendigen Ressourcen frühzeitig sichern und die Implementierung parallel zur Projektvorbereitung vorantreiben, um von Beginn an eine wirksame Kostensteuerung aufrechtzuerhalten.